\def\tituloconcepto{Medición y regla de Born}
% Si el capitulo se compila solo (sin main.tex), esta cabecera abre el documento.
% Antes de % Si el capitulo se compila solo (sin main.tex), esta cabecera abre el documento.
% Antes de % Si el capitulo se compila solo (sin main.tex), esta cabecera abre el documento.
% Antes de \input{cabecera} cada capitulo define \tituloconcepto.
\ifdefined\enmain\else
  \documentclass[11pt,a4paper]{article}
  \input{preambulo}
  \begin{document}
  \title{\tituloconcepto}
  \author{\autor}
  \date{\today}
  \maketitle
\fi
 cada capitulo define \tituloconcepto.
\ifdefined\enmain\else
  \documentclass[11pt,a4paper]{article}
  % Preambulo compartido por todos los documentos. Compilar con XeLaTeX/LuaLaTeX
% (o Tectonic). Si se usa Overleaf, elegir el compilador XeLaTeX.
\usepackage[spanish,es-noshorthands]{babel}
\usepackage{amsmath,amssymb}
\usepackage{booktabs,array}
\usepackage[a4paper,margin=2.5cm]{geometry}
\usepackage{xcolor}
\usepackage{listings}
\usepackage[most]{tcolorbox}
\usepackage{quantikz}
\usepackage[hidelinks]{hyperref}

% ---- Autor (editar a gusto) -------------------------------------------------
\newcommand{\autor}{Genesis Silesky}

% ---- Notacion de Dirac y atajos ---------------------------------------------
% quantikz ya define \ket y \bra; DeclareRobustCommand los redefine sin error.
\DeclareRobustCommand{\ket}[1]{\left|#1\right\rangle}
\DeclareRobustCommand{\bra}[1]{\left\langle#1\right|}
\DeclareRobustCommand{\braket}[2]{\left\langle#1\middle|#2\right\rangle}
\newcommand{\C}{\mathbb{C}}
\newcommand{\R}{\mathbb{R}}
\DeclareMathOperator{\Tr}{Tr}

% ---- Cajas de apoyo --------------------------------------------------------
% "concepto": definiciones previas que hacen falta para entender el capitulo.
\newtcolorbox{concepto}[1]{breakable,colback=blue!4,colframe=blue!55!black,
  fonttitle=\bfseries,title={Concepto previo: #1}}
% "nota": aclaraciones y advertencias.
\newtcolorbox{nota}[1]{breakable,colback=orange!6,colframe=orange!70!black,
  fonttitle=\bfseries,title={#1}}

% ---- Estilos de listados ----------------------------------------------------
\lstdefinestyle{py}{language=Python,basicstyle=\ttfamily\footnotesize,
  breaklines=true,postbreak=\mbox{\textcolor{gray}{$\hookrightarrow$}\space},
  numbers=left,numberstyle=\tiny\color{gray},frame=single,
  commentstyle=\color{green!40!black}\itshape,
  keywordstyle=\color{blue!70!black}\bfseries,
  stringstyle=\color{red!60!black},showstringspaces=false,tabsize=4,
  columns=fullflexible,keepspaces=true}
\lstdefinestyle{salida}{basicstyle=\ttfamily\footnotesize,breaklines=true,
  frame=leftline,backgroundcolor=\color{gray!8},columns=fullflexible,
  keepspaces=true}

  \begin{document}
  \title{\tituloconcepto}
  \author{\autor}
  \date{\today}
  \maketitle
\fi
 cada capitulo define \tituloconcepto.
\ifdefined\enmain\else
  \documentclass[11pt,a4paper]{article}
  % Preambulo compartido por todos los documentos. Compilar con XeLaTeX/LuaLaTeX
% (o Tectonic). Si se usa Overleaf, elegir el compilador XeLaTeX.
\usepackage[spanish,es-noshorthands]{babel}
\usepackage{amsmath,amssymb}
\usepackage{booktabs,array}
\usepackage[a4paper,margin=2.5cm]{geometry}
\usepackage{xcolor}
\usepackage{listings}
\usepackage[most]{tcolorbox}
\usepackage{quantikz}
\usepackage[hidelinks]{hyperref}

% ---- Autor (editar a gusto) -------------------------------------------------
\newcommand{\autor}{Genesis Silesky}

% ---- Notacion de Dirac y atajos ---------------------------------------------
% quantikz ya define \ket y \bra; DeclareRobustCommand los redefine sin error.
\DeclareRobustCommand{\ket}[1]{\left|#1\right\rangle}
\DeclareRobustCommand{\bra}[1]{\left\langle#1\right|}
\DeclareRobustCommand{\braket}[2]{\left\langle#1\middle|#2\right\rangle}
\newcommand{\C}{\mathbb{C}}
\newcommand{\R}{\mathbb{R}}
\DeclareMathOperator{\Tr}{Tr}

% ---- Cajas de apoyo --------------------------------------------------------
% "concepto": definiciones previas que hacen falta para entender el capitulo.
\newtcolorbox{concepto}[1]{breakable,colback=blue!4,colframe=blue!55!black,
  fonttitle=\bfseries,title={Concepto previo: #1}}
% "nota": aclaraciones y advertencias.
\newtcolorbox{nota}[1]{breakable,colback=orange!6,colframe=orange!70!black,
  fonttitle=\bfseries,title={#1}}

% ---- Estilos de listados ----------------------------------------------------
\lstdefinestyle{py}{language=Python,basicstyle=\ttfamily\footnotesize,
  breaklines=true,postbreak=\mbox{\textcolor{gray}{$\hookrightarrow$}\space},
  numbers=left,numberstyle=\tiny\color{gray},frame=single,
  commentstyle=\color{green!40!black}\itshape,
  keywordstyle=\color{blue!70!black}\bfseries,
  stringstyle=\color{red!60!black},showstringspaces=false,tabsize=4,
  columns=fullflexible,keepspaces=true}
\lstdefinestyle{salida}{basicstyle=\ttfamily\footnotesize,breaklines=true,
  frame=leftline,backgroundcolor=\color{gray!8},columns=fullflexible,
  keepspaces=true}

  \begin{document}
  \title{\tituloconcepto}
  \author{\autor}
  \date{\today}
  \maketitle
\fi


\section{Medición y regla de Born}

\subsection{Objetivo}
Comprobar experimentalmente que $P(0)=|\alpha|^2$ y $P(1)=|\beta|^2$, y ver
cómo la frecuencia observada converge a esas probabilidades al aumentar los
\emph{shots}.

\subsection{Conceptos previos}
\begin{concepto}{regla de Born, colapso y frecuencia relativa}
Ver sección~\ref{sec:prerreq}. Dos ideas: (i) medir produce un resultado
clásico aleatorio con las probabilidades de Born; (ii) tras medir, el estado
\emph{colapsa} al resultado observado, y la superposición original se pierde.
\end{concepto}

\subsection{Fundamento teórico}
Para $\ket\psi=\alpha\ket0+\beta\ket1$, la medición en la base computacional
da $0$ con probabilidad $|\alpha|^2$ y $1$ con $|\beta|^2$. No se puede
conocer $\alpha,\beta$ con una sola medición; hay que repetir el
experimento $N$ veces (cada repetición es un \emph{shot}). La frecuencia
$\hat p$ de ``1'' tiene error típico
\begin{equation}
  \sigma=\sqrt{\frac{p(1-p)}{N}} .
\end{equation}
Para $p=0.25$: $\sigma\approx0.137$ ($N=10$), $0.043$ ($N=100$),
$0.014$ ($N=1000$), $0.0043$ ($N=10000$).

\paragraph{Preparación del estado.} La rotación
$R_y(\theta)\ket0=\cos\frac\theta2\ket0+\sin\frac\theta2\ket1$. Para que
$P(1)=\sin^2\frac\theta2=0.25$ se necesita $\theta=2\arcsin(0.5)=\pi/3$.

\begin{center}
\begin{quantikz}
\lstick{$q_0:\ \ket0$} & \gate{R_y(\pi/3)} & \meter{}
\end{quantikz}
\end{center}

\subsection{Código (completamente comentado)}
\lstinputlisting[style=py]{../04_medicion.py}

\subsection{Explicación celda por celda}
\begin{description}
  \item[Celda 1] Además de las importaciones, reconfigura la salida estándar a
    UTF-8: sin eso la consola de Windows no puede imprimir el dibujo del
    circuito (error \texttt{UnicodeEncodeError}).
  \item[Celda 2] Calcula $\theta$, crea un circuito con $1$ qubit y $1$ bit
    clásico, aplica $R_y$ y mide.
  \item[Celda 3] Ejecuta $10,100,1000,10000$ shots y calcula la frecuencia de
    ``1''. \texttt{counts.get('1', 0)} evita un error si ``1'' nunca salió.
  \item[Celda 4] Dibuja el histograma (en script puede guardarse con
    \texttt{savefig}).
\end{description}

\subsection{Salida obtenida}
\begin{nota}{Los números cambian en cada ejecución}
La medición es aleatoria: los conteos exactos varían. Lo que debe
reproducirse es la \emph{tendencia}: la frecuencia de ``1'' se acerca a
$0.25$ al aumentar los shots.
\end{nota}
\lstinputlisting[style=salida]{../salidas/04_medicion.txt}

\subsection{Interpretación}
Con $10$ shots la frecuencia puede alejarse mucho de $0.25$ (error típico
$\approx0.14$); con $10\,000$ queda cerca (error $\approx0.004$). La
incertidumbre estadística disminuye como $1/\sqrt N$.

\subsection{Ejercicios}
\begin{enumerate}
  \item Prepare un qubit con $P(1)=0.9$ y mida $2000$ shots.
  \item Repita 20 veces la corrida de 100 shots y grafique la dispersión de
    la frecuencia; compárela con $\sigma=\sqrt{p(1-p)/N}$.
  \item ¿Qué obtiene si mide dos veces seguidas el mismo qubit? (Pista:
    colapso.)
\end{enumerate}

\ifdefined\enmain\else
  \end{document}
\fi

